% GENERATED FILE. Edit canonical lessons, then run npm run generate:books.
% canonical-source-hash: fnv1a64:9353148ca2b32d1f
% canonical-lessons: JA-C41-tegami, JA-C41-jikan, JA-C41-namae, JA-C41-okane, JA-C41-yama

\chapter{Letter, Time, Name, Money, Mountain}
\label{ch:ja-letter-time-name-money-mountain}
\hlchaptermodality{41}

\begin{chapteropening}
\textbf{By the end of this chapter:} \emph{I can say a letter, time, a name, money and a mountain.}

Complete the last lesson of chapter 41: I can say a letter, time, a name, money and a mountain.
\end{chapteropening}

\section[\texorpdfstring{\ja{てがみ} (tegami)}{てがみ (tegami)}]{\ja{てがみ} (tegami) — a letter (post)}

\label{lesson:JA-C41-tegami}



\begin{quote}
Before the new one: say the Japanese for a problem, then the Japanese for an answer.
\end{quote}

\subsection*{You'll want to know: \ja{てがみ}}
\textbf{\ja{てがみ}} — \emph{tegami} — \textquotedblleft{}a letter (post)\textquotedblright{}.

\begin{cousinweb}
\textbf{\ja{てがみ}} (\emph{tegami}) is a compound: native \emph{te}, \textquotedblleft{}hand\textquotedblright{}, and \emph{kami}, \textquotedblleft{}paper\textquotedblright{}, voiced to \emph{gami}, both kun'yomi of its kanji (JMdict, KANJIDIC2). \emph{Kami} is usually counted native, but an early loan from Chinese has also been suggested; its origin is not settled.
\end{cousinweb}

\begin{grammarlens}[title={What you've built}]
One more word for this chapter.
\end{grammarlens}

\subsection*{Your turn}
\begin{itemize}
\raggedright
  \item \emph{Say it:} \emph{tegami}
  \item \emph{Say it:} \emph{tegami}, once more
  \item \emph{Say it:} say \emph{tegami}
  \item \emph{Recall:} say the Japanese for a problem, then the Japanese for an answer, then say \emph{tegami} again
\end{itemize}

\begin{tcolorbox}[breakable,colback=teal!4,colframe=teal!35!black,title={Before you move on}]
What does \ja{てがみ} mean? (\textquotedblleft{}A letter (post)\textquotedblright{}.) Say it once more.
\end{tcolorbox}
\section[\texorpdfstring{\ja{じかん} (jikan)}{じかん (jikan)}]{\ja{じかん} (jikan) — time; an hour}

\label{lesson:JA-C41-jikan}



\begin{quote}
Before the new one: say the Japanese for an answer, then the Japanese for a letter.
\end{quote}

\subsection*{You'll want to know: \ja{じかん}}
\textbf{\ja{じかん}} — \emph{jikan} — \textquotedblleft{}time; an hour\textquotedblright{}.

\begin{cousinweb}
\textbf{\ja{じかん}} (\emph{jikan}) is Sino-Japanese (kango): two Chinese characters read with their on'yomi, readings borrowed from Chinese. They are \emph{ji}, \textquotedblleft{}time\textquotedblright{} (Middle Chinese \emph{dzyi}) and \emph{kan}, \textquotedblleft{}interval\textquotedblright{} (\emph{kean}) (Unihan, KANJIDIC2; Middle Chinese in Baxter's transcription).
\end{cousinweb}

\begin{grammarlens}[title={What you've built}]
One more word for this chapter.
\end{grammarlens}

\subsection*{Your turn}
\begin{itemize}
\raggedright
  \item \emph{Say it:} \emph{jikan}
  \item \emph{Say it:} \emph{jikan}, once more
  \item \emph{Say it:} say \emph{jikan}
  \item \emph{Recall:} say the Japanese for an answer, then the Japanese for a letter, then say \emph{jikan} again
\end{itemize}

\begin{tcolorbox}[breakable,colback=teal!4,colframe=teal!35!black,title={Before you move on}]
What does \ja{じかん} mean? (\textquotedblleft{}Time; an hour\textquotedblright{}.) Say it once more.
\end{tcolorbox}
\section[\texorpdfstring{\ja{なまえ} (namae)}{なまえ (namae)}]{\ja{なまえ} (namae) — a name}

\label{lesson:JA-C41-namae}



\begin{quote}
Before the new one: say the Japanese for a letter, then the Japanese for time; an hour.
\end{quote}

\subsection*{You'll want to know: \ja{なまえ}}
\textbf{\ja{なまえ}} — \emph{namae} — \textquotedblleft{}a name\textquotedblright{}.

\begin{cousinweb}
\textbf{\ja{なまえ}} (\emph{namae}) is a native compound, inherited rather than borrowed: \emph{na}, \textquotedblleft{}name\textquotedblright{}, and \emph{mae}, \textquotedblleft{}front\textquotedblright{}, both kun'yomi of its kanji (JMdict, KANJIDIC2).
\end{cousinweb}

\begin{grammarlens}[title={What you've built}]
One more word for this chapter.
\end{grammarlens}

\subsection*{Your turn}
\begin{itemize}
\raggedright
  \item \emph{Say it:} \emph{namae}
  \item \emph{Say it:} \emph{namae}, once more
  \item \emph{Say it:} say \emph{namae}
  \item \emph{Recall:} say the Japanese for a letter, then the Japanese for time; an hour, then say \emph{namae} again
\end{itemize}

\begin{tcolorbox}[breakable,colback=teal!4,colframe=teal!35!black,title={Before you move on}]
What does \ja{なまえ} mean? (\textquotedblleft{}A name\textquotedblright{}.) Say it once more.
\end{tcolorbox}
\section[\texorpdfstring{\ja{おかね} (okane)}{おかね (okane)}]{\ja{おかね} (okane) — money}

\label{lesson:JA-C41-okane}



\begin{quote}
Before the new one: say the Japanese for time; an hour, then the Japanese for a name.
\end{quote}

\subsection*{You'll want to know: \ja{おかね}}
\textbf{\ja{おかね}} — \emph{okane} — \textquotedblleft{}money\textquotedblright{}.

\begin{cousinweb}
\textbf{\ja{おかね}} (\emph{okane}) is the polite prefix \emph{o-} with native \emph{kane}, \textquotedblleft{}metal; money\textquotedblright{}, read as the kun'yomi of the character for gold and metal. It is inherited rather than borrowed (JMdict, KANJIDIC2).
\end{cousinweb}

\begin{grammarlens}[title={What you've built}]
One more word for this chapter.
\end{grammarlens}

\subsection*{Your turn}
\begin{itemize}
\raggedright
  \item \emph{Say it:} \emph{okane}
  \item \emph{Say it:} \emph{okane}, once more
  \item \emph{Say it:} say \emph{okane}
  \item \emph{Recall:} say the Japanese for time; an hour, then the Japanese for a name, then say \emph{okane} again
\end{itemize}

\begin{tcolorbox}[breakable,colback=teal!4,colframe=teal!35!black,title={Before you move on}]
What does \ja{おかね} mean? (\textquotedblleft{}Money\textquotedblright{}.) Say it once more.
\end{tcolorbox}
\section[\texorpdfstring{\ja{やま} (yama)}{やま (yama)}]{\ja{やま} (yama) — a mountain}

\label{lesson:JA-C41-yama}



\begin{quote}
Before the new one: say the Japanese for a name, then the Japanese for money.
\end{quote}

\subsection*{You'll want to know: \ja{やま}}
\textbf{\ja{やま}} — \emph{yama} — \textquotedblleft{}a mountain\textquotedblright{}.

\begin{cousinweb}
\textbf{\ja{やま}} (\emph{yama}) is a native Japanese word, inherited rather than borrowed from Chinese, and its older origin is not settled. Its kanji was attached to it later as a kun'yomi, a native reading given to a Chinese character (JMdict, KANJIDIC2).
\end{cousinweb}

\begin{grammarlens}[title={What you've built}]
One more word for this chapter.
\end{grammarlens}

\subsection*{Your turn}
\begin{itemize}
\raggedright
  \item \emph{Say it:} \emph{yama}
  \item \emph{Say it:} \emph{yama}, once more
  \item \emph{Say it:} say \emph{yama}
  \item \emph{Recall:} say the Japanese for a name, then the Japanese for money, then say \emph{yama} again
\end{itemize}

\begin{tcolorbox}[breakable,colback=teal!4,colframe=teal!35!black,title={Before you move on}]
What does \ja{やま} mean? (\textquotedblleft{}A mountain\textquotedblright{}.) Say it once more.
\end{tcolorbox}
